% Coupling paper, first complete draft (27 Oct 2026).
% Build (offline): pdflatex main; bibtex main; pdflatex main; pdflatex main.
% The Springer Nature template (sn-jnl.cls, Biomechanics and Modeling in
% Mechanobiology) is not installed here; the manuscript uses the article class
% with natbib (author-year), 12 pt, single column and line numbers. The Springer
% template will be applied at submission.
\documentclass[12pt,a4paper]{article}
\usepackage[utf8]{inputenc}
\usepackage[T1]{fontenc}
\usepackage{lmodern}
\usepackage[a4paper,margin=25mm]{geometry}
\usepackage{amsmath,amssymb,bm,mathtools}
\usepackage{graphicx}
\usepackage{booktabs}
\usepackage{array}
\usepackage{caption}
\usepackage{subcaption}
\usepackage[authoryear,round]{natbib}
\usepackage{algorithm}
\usepackage{algpseudocode}
\usepackage{lineno}
\usepackage{xcolor}
\usepackage[colorinlistoftodos,textsize=footnotesize]{todonotes}
\usepackage[hidelinks]{hyperref}
\usepackage{microtype}
\setlength{\emergencystretch}{2em}
\captionsetup{font=small,labelfont=bf}
\linenumbers

% Paths of the figures (not copied into this folder).
\newcommand{\assets}{../../../assets/}
\newcommand{\figs}{../}
% Figure in preparation by another script: box with the file name as a fallback.
\newcommand{\figorbox}[2]{\IfFileExists{#1}{\includegraphics[#2]{#1}}{\fbox{\parbox[c][5cm][c]{0.9\linewidth}{\centering figure file \texttt{\detokenize{#1}} not yet generated}}}}

% Notation (Klempt et al. 2024, 2026).
\newcommand{\phiD}{\phi_{3D}}
\newcommand{\lv}{\textstyle\sum_i\bar\phi_i}
\newcommand{\nc}{\mathbf n_c}
\newcommand{\fterm}{\ensuremath{-r\tfrac{c}{k+c}\,\nc\cdot\nabla\phiD}}
\newcommand{\TODO}[1]{\todo[inline]{TODO: #1}}

\title{Coupling a multi-species point model to a continuum biofilm growth model: verification in three codes and the effect of the composition on growth and stress}
\author{Authors to be agreed\thanks{Author list, order and affiliations to be agreed by all contributors before the manuscript is shared; corresponding author to be named.}}
\date{Draft of 27 October 2026}

\begin{document}
\maketitle

\begin{abstract}
Continuum models of biofilm growth predict where biomass forms and the stresses it causes, but treat the biofilm as one material. Multi-species point models describe the composition, the living fraction and the response to an antibiotic, but have no space. This work couples the multi-species point model of Klempt et al.\ (2026), used as published, to the Hamilton-principle growth model of Klempt et al.\ (2024) at every integration point, without changing the equations of either model. The point model receives the local nutrient and antibiotic from the fields, and its living volume fraction $\sum_i\bar\phi_i$ sets the consumption $g\,\phi_{3D}(\sum_i\bar\phi_i)\,c$ of both. The coupling is implemented in a finite-volume code, in a neighboured-element user element in ANSYS and in a Galerkin user element in Abaqus, and verified by replay of the point model, a manufactured solution and grid convergence indices. The advection term of the growth equation has a cell P\'eclet number of about 31 on a $16^3$ grid; a constant added diffusion gives a quarter of the converged seed stress, a second-order treatment is needed. With it the growth converges in all examples (GCI 0.7 to 3.3\,\%) and the seed stress converges when the species start apart (GCI 0.7 to 8\,\%). An antibiotic removes most of the living biomass of the seed (0.97 to 0.15 when the species start apart, 0.96 to 0.04 when mixed), while the growth changes by about 1\,\% and the seed stress by about 10\,\%, because the growth law counts all biomass. The local nutrient decides where the antibiotic acts. The composition that newly reached material receives is not fixed by either model and decides, in a maze, whether the biofilm recovers after the antibiotic.
\end{abstract}

\noindent\textbf{Keywords:} biofilm growth; multi-species model; antibiotic; extended Hamilton principle; code verification; grid convergence index; growth-induced stress

\section{Introduction}
\label{sec:intro}

Oral biofilms on teeth and implants are communities of several bacterial species. Their composition changes over days: early colonisers prepare the surface for later ones, and pathogenic species appear only when the community allows it \citep{Heine2025PeriImplant}. Treatments such as antibiotics act on the composition, while the mechanical load on the surface comes from the amount and the growth of the biomass. A model that predicts the stresses for a given treatment therefore has to describe both: where the biofilm grows, and what it is made of.

Two kinds of model exist for this. Continuum growth models describe the amount of biomass in space. \citet{WannerGujer1986} introduced the one-dimensional multispecies biofilm, \citet{Eberl2001Biofilm} a density-dependent diffusion model of the spreading. \citet{Feng2021SymbioticBiofilm} model a two-species oral biofilm in which the spreading follows the summed growth of the species, and \citet{Soleimani2023CoAggregation} the co-aggregation of two species in a user element. \citet{Klempt2024DiffusionDrivenGrowth} derive the growth of a biofilm from the extended Hamilton principle \citep{JunkerBalzani2021ExtendedHamilton}, with a phase field $\phiD$, a nutrient $c$ and the growth variable $\alpha$ of a multiplicative split $\mathbf F=\mathbf F_e\mathbf F_g$, and compute the resulting stresses; their model has one biomass and no composition. Point models describe the composition of a community at one point. \citet{Klempt2026ContinuumBacterialGrowth} derive a multi-species model from the same variational principle, with volume fractions $\phi_i$ and living fractions $\psi_i$ of $N$ species, an interaction matrix $\mathbf A$ and the response of each species to an antibiotic; \citet{Fritsch2025BayesianMicrofilms} calibrate its parameters by Bayesian updating. These models have no space: the nutrient $c^*$ and the antibiotic $\alpha^*$ are given inputs.

The point model needs local values of $c^*$ and $\alpha^*$, which a growth field can supply. The growth field needs to know how much of the biomass is alive, because only living biomass consumes nutrient. Connecting the two at every material point is the natural step, but it raises questions that neither paper answers: which variables are passed, whether either model has to be changed, which new assumptions are needed (for example the composition of biofilm that a point reaches later), and whether the composition reaches the growth and the stress at all.

This work couples the point model of \citet{Klempt2026ContinuumBacterialGrowth} to the growth model of \citet{Klempt2024DiffusionDrivenGrowth} without changing either set of equations. The contribution is threefold.
\begin{enumerate}
\item The point model runs at every integration point as published. It receives $c^*=c^*_0c$ and $\alpha^*=\alpha^*_0a$ from the nutrient and antibiotic fields and returns the living volume fraction $\lv$, which sets the consumption $g\,\phiD(\lv)\,c$ in the field equations.
\item The coupling is implemented three times, in a finite-volume code, in a neighboured-element user element in ANSYS and in a Galerkin user element in Abaqus, and verified: replay of the point model alone, a manufactured solution, agreement of the codes and grid convergence with the grid convergence index. The advection term \fterm{} of the growth equation turns out to set the accuracy of the stress.
\item Examples with two species and an antibiotic, with the synthetic data of \citet{Fritsch2025BayesianMicrofilms} and the constants of \citet{Klempt2024DiffusionDrivenGrowth,Klempt2026ContinuumBacterialGrowth}, show what the composition changes and what it does not, and which assumption beyond the two papers matters.
\end{enumerate}
Section~\ref{sec:models} states the models and the coupling, Section~\ref{sec:verification} the implementations and their verification, Section~\ref{sec:results} the examples, Section~\ref{sec:discussion} the limits of the coupling, and Section~\ref{sec:conclusions} concludes.

\section{Models and coupling}
\label{sec:models}

\subsection{Growth field}
\label{sec:field}

The biofilm is described by the phase field $\phiD\in[0,1]$, the nutrient $c$ and the growth variable $\alpha$ of \citet{Klempt2024DiffusionDrivenGrowth}, with $\mathbf F=\mathbf F_e\mathbf F_g$, $\mathbf F_g=\alpha\mathbf I$ and $\alpha(0)=1$, so that the growth is $\alpha-1$. After division by the viscosities, their Eqs.~34 to 36 read
\begin{align}
\dot\phiD &= \beta\,\nabla^2\phiD + k_\alpha\,\alpha - r\,\frac{c}{k+c}\,\nc\cdot\nabla\phiD, \qquad \nc=\frac{\nabla c}{|\nabla c|}, \label{eq:34}\\
\dot c &= d\,\nabla^2 c - g\,\phiD\Bigl(\lv\Bigr)\,c, \label{eq:35}\\
\dot\alpha &= k_\alpha\,\phiD, \label{eq:36}
\end{align}
with zero flux on the boundary except where the nutrient is held. The last term of Eq.~\eqref{eq:34} moves the edge of the biofilm up the nutrient gradient with the speed $r\,c/(k+c)$; we call it the front term. Two details differ from the equations as printed, both confirmed by the authors: the consumption of Eq.~\eqref{eq:35} is $g\,\phi\,c$ (the printed $g\,\phi$ is a misprint), and the front direction is regularised as $\nabla c/(|\nabla c|^2+\varepsilon^2)^{1/2}$ with $\varepsilon^2=10^{-6}$ ($\varepsilon^2=10^{-8}$ changes the seed stress by $0.1\,\%$). The factor $\lv$ in Eq.~\eqref{eq:35} is the coupling term of this work (Section~\ref{sec:coupling}); in \citet{Klempt2024DiffusionDrivenGrowth} the whole biofilm consumes. The mechanical term $\phi\,\mu\,(\mathbf I:\mathbf C_e-3)$ of their Eq.~30, which they neglect because $\mu$ is small against the other parameters, is neglected here as well, so the mechanics is coupled one way. The elastic response is neo-Hookean with the stiffness $E(\phiD)=E\,(\phiD^2+f)$, $E=10$~Pa, $\nu=0.49$ and a void floor $f=10^{-3}$ (an assumption that keeps the empty cells solvable; $f\to0$ changes the seed stress by $0.13\,\%$).

\subsection{Point model}
\label{sec:point}

At a material point the multi-species model of \citet{Klempt2026ContinuumBacterialGrowth} describes the volume fractions $\phi_i$ and the living fractions $\psi_i$ of $N$ species and the void $\phi_0$, with $\sum_i\phi_i+\phi_0=1$. It follows from the extended Hamilton principle with the free energy $\Psi=-\tfrac12c^*\,\bar{\bm\phi}^{\mathrm T}\mathbf A\,\bar{\bm\phi}$, $\bar\phi_i=\phi_i\psi_i$, a dissipation with the viscosities $\eta_i$, $\eta_{\phi,i}$ and the constraint enforced by the multiplier $\gamma$:
\begin{equation}
\begin{aligned}
(\eta_{\phi,i}+\eta_i\psi_i^2)\,\dot\phi_i+\eta_i\phi_i\psi_i\,\dot\psi_i
&= c^*\psi_i\textstyle\sum_j A_{ij}\bar\phi_j-\gamma-(\text{barrier}),\\
\eta_i\,(\phi_i\psi_i\,\dot\phi_i+\phi_i^2\,\dot\psi_i)
&= c^*\phi_i\textstyle\sum_j A_{ij}\bar\phi_j-\alpha^* b_i\,\psi_i-(\text{barrier}),
\end{aligned}
\label{eq:point}
\end{equation}
where $A_{ij}$ are the interactions, $b_i$ the sensitivity of species $i$ to an antibiotic of concentration $\alpha^*$, $c^*$ the nutrient, and the barrier terms keep $\phi_i,\psi_i$ in $(0,1)$. The antibiotic lowers $\psi_i$; the point model compares $c^*$ with $\alpha^*b_i$. Our implementation reproduces the fourteen numerical examples of the paper (two printed details read differently: Eq.~17 as printed carries a $+\gamma$ that the figures do not show, and the four-species interaction matrix of Eq.~19 reproduces the figures without its prefactor $\tfrac12$) and the figures of \citet{Fritsch2025BayesianMicrofilms} within $0.013$. The model is used as published, with its own constraint and its implicit time integration (steps of $10^{-4}$ in its own time unit, Newton's method). For the examples the synthetic two-species data of \citet{Fritsch2025BayesianMicrofilms}, case~I, are used with their true parameters: $\mathbf A=\begin{psmallmatrix}1&0.1\\0.1&1\end{psmallmatrix}$, $b=(1,2)$, $\eta=(1,2)$, $c^*_0=100$.

\subsection{Coupling}
\label{sec:coupling}

The two models are connected at every integration point without changing either (Figure~\ref{fig:schematic}, Algorithm~\ref{alg:coupled}).
\begin{enumerate}
\item \emph{Fields to point model.} The point model sees $c^*=c^*_0\,c$ and $\alpha^*=\alpha^*_0\,a$, with $a$ the antibiotic, transported by a second field with the diffusivity and uptake of the nutrient, $\dot a=d\,\nabla^2a-g\,\phiD(\lv)\,a$. Between the increments $t_n$ and $t_{n+1}$ of the field both are interpolated linearly in time over the inner steps of the point model,
\begin{equation}
c^*(t)=c^*_0\Bigl[\tfrac{t_{n+1}-t}{\Delta t}\,c_n+\tfrac{t-t_n}{\Delta t}\,c_{n+1}\Bigr],\qquad t_n\le t\le t_{n+1}.
\label{eq:cstar}
\end{equation}
\citet{Klempt2026ContinuumBacterialGrowth} already use a time-dependent $c^*$ (their four-species case~3), so the point model accepts such an input without change.
\item \emph{Point model to fields.} Only the living biomass consumes: the uptake of nutrient and antibiotic is $g\,\phiD(\lv)$, Eq.~\eqref{eq:35}. The point model normalises its volume, so $\bar\phi_i\in[0,1]$ is the living volume fraction of species $i$ at the point and $\phiD\lv\le\phiD$ the living volume fraction in the field; the void $\phi_0$ has no living fraction and is not part of the sum. A dead biofilm consumes nothing. The growth law~\eqref{eq:36} is unchanged: $\alpha$ measures the biomass produced, alive or dead.
\item \emph{Time.} \citet{Klempt2024DiffusionDrivenGrowth} normalise time with a reference time chosen for the problem at hand. We take $T^*=1$ as the duration covered by the point model's data, so that both models run on the same time: the point model advances $s\,\Delta T^*$ per increment, with $s$ the length of that window in its own units, not a fitted value ($s=0.1$ for case~I, 1000 steps of $10^{-4}$).
\item \emph{Start of the point model.} It starts at a point once $\phiD$ there reaches $\phi_{\min}=0.001$. Where a seed is placed it starts from the seed composition. A point that the spreading biofilm reaches later takes the living composition of its colonised neighbours (their amounts divided in the ratio of $\bar\phi_i$) and starts alive ($\psi_i=0.999$, $\gamma=0$), because the biofilm spreads through the growth of its living cells. The papers do not fix this rule; it is an assumption, and the alternative of starting from the seed composition is compared in Section~\ref{sec:newpoints}.
\end{enumerate}
The composition reaches the growth and the stress only through the nutrient field and the front term of Eq.~\eqref{eq:34}: the antibiotic lowers $\psi_i$, less nutrient is consumed, more reaches the front, and the edge of the biofilm moves differently.

\begin{figure}[tbp]
  \centering
  \includegraphics[width=\linewidth]{\assets fig_coupling_schematic_agreed.png}
  \caption{The coupling: the fields $\phiD$, $c$ and $a$ of the growth model, the point model at every integration point with the inputs $c^*=c^*_0c$, $\alpha^*=\alpha^*_0a$ and the output $\lv$, and the mechanics with $\mathbf F_g=\alpha\mathbf I$.}
  \label{fig:schematic}
\end{figure}

\begin{algorithm}[tbp]
\caption{One increment $t_n\to t_{n+1}$ of the coupled run (user-element form)}
\label{alg:coupled}
\begin{algorithmic}[1]
\Require at every integration point $\alpha_n$, the point-model state $(\phi_i,\psi_i,\gamma)_n$ and the pair $(c,a)$ seen in the previous increment; the fields $\phiD$, $c$, $a$ of $t_n$
\Ensure displacements, stresses and state of $t_{n+1}$
\Repeat \Comment{Newton iteration of the equilibrium}
  \ForAll{integration points}
    \State $\alpha_{n+1}\leftarrow\alpha_n+k_\alpha\,\phiD\,\Delta t$ \Comment{Eq.~\eqref{eq:36}}
    \State $\mathbf F_e\leftarrow\mathbf F\,\mathbf F_g^{-1}$, $\mathbf F_g=\alpha_{n+1}\mathbf I$; stress and tangent from $W(\mathbf F_e)$ and $E(\phiD)$
    \If{$\phiD\ge\phi_{\min}$ or the point model has started}
      \If{not yet advanced in this increment}
        \State $(c_n,c_{n+1})$, $(a_n,a_{n+1})$ from the previous and the current value at the point
        \State $n_{\rm sub}$ steps of $s\Delta t/n_{\rm sub}$ of Eq.~\eqref{eq:point} with $c^*$, $\alpha^*$ from Eq.~\eqref{eq:cstar}
        \State store $(\phi_i,\psi_i,\gamma)_{n+1}$ and $\lv$
      \EndIf
    \EndIf
  \EndFor
\Until{the residual meets the convergence criterion}
\State advance $\phiD$ (Eq.~\eqref{eq:34}), $c$ and $a$ (Eq.~\eqref{eq:35} with $\lv$)
\end{algorithmic}
\end{algorithm}

\subsection{Parameters}
\label{sec:parameters}

All constants are those of the papers (Table~\ref{tab:parameters}): Table~2 of \citet{Klempt2024DiffusionDrivenGrowth}, converted from their $20~\mu$m cube to the 2~mm cube of the examples, and the antibiotic of case~4 of \citet{Klempt2026ContinuumBacterialGrowth}, $\alpha^*_0=100$ from $T^*=1/3$. With $d=10^8$~mm$^2/T^*$ the fields adjust within about $4\times10^{-8}\,T^*$, so the fields are quasi-steady; keeping the time derivative of Eq.~\eqref{eq:35} changes the seed values by less than $0.005$. A diffusion slower by eight orders of magnitude ($d=1$~mm$^2/T^*$) is only a sensitivity: the antibiotic then reaches the seed more slowly and kills markedly less of it. Values not given by the papers are marked as assumptions: the length and time scale of the conversion, $\phi_{\min}$ (results change by less than $1\,\%$ below $0.001$), the start rule of new points, the void floor $f$ and the regularisation $\varepsilon$.

\begin{table}[tbp]
  \centering
  \caption{Constants of the examples. Conversion: Table~2 of \citet{Klempt2024DiffusionDrivenGrowth} on the 2~mm cube with the half edge $\ell=1$~mm, so that $g\ell^2/d=1$.}
  \label{tab:parameters}
  \footnotesize
  \setlength{\tabcolsep}{4pt}
  \begin{tabular}{lll>{\raggedright\arraybackslash}p{5.2cm}}
    \toprule
    quantity & symbol & value & source \\
    \midrule
    diffusion of $\phiD$ & $\beta$ & $0.02$~mm$^2/T^*$ & \citet{Klempt2024DiffusionDrivenGrowth}, converted (assumption) \\
    front speed, half saturation & $r$, $k$ & $10$~mm$/T^*$, $1$ & \citet{Klempt2024DiffusionDrivenGrowth}, converted \\
    growth rate & $k_\alpha$ & $10^{-3}\,T^{*-1}$ & \citet{Klempt2024DiffusionDrivenGrowth} \\
    nutrient diffusion, uptake & $d$, $g$ & $10^8$~mm$^2/T^*$, $10^8\,T^{*-1}$ & \citet{Klempt2024DiffusionDrivenGrowth}, converted \\
    antibiotic diffusion, uptake & $d_a$, $g_a$ & $=d$, $=g$ & assumption \\
    stiffness, Poisson ratio & $E$, $\nu$ & $10$~Pa, $0.49$ & \citet{Klempt2024DiffusionDrivenGrowth} \\
    void floor & $f$ & $10^{-3}$ & assumption \\
    point model & $\mathbf A$, $b$, $\eta$, $c^*_0$ & case~I & \citet{Fritsch2025BayesianMicrofilms} \\
    antibiotic level, onset & $\alpha^*_0$, $T^*_{\rm ab}$ & $100$, $1/3$ & \citet{Klempt2026ContinuumBacterialGrowth}, case~4 \\
    time link & $s$ & $0.1$ & from the data of case~I \\
    start threshold & $\phi_{\min}$ & $0.001$ & assumption (converged) \\
    front regularisation & $\varepsilon^2$ & $10^{-6}$ & assumption \\
    \bottomrule
  \end{tabular}
\end{table}

\section{Implementations and verification}
\label{sec:verification}

\subsection{Three codes}

The coupled model is solved in three codes that share the equations, the parameters and the compiled point-model routine, but not the discretisation (Table~\ref{tab:codes}). The ANSYS element is the neighboured element method of the original authors' group \citep{Blaszczyk2022TaylorWLS,Rudolf2025NEM}: the Laplacian and the gradients are fitted on the Gauss points from 30 neighbours by a weighted least-squares Taylor expansion, and the fields are updated explicitly. Three numerical details of this element had to be corrected before the results below, without changing its equations (Appendix~\ref{app:element}). Differences between the codes therefore come from the fields and the mechanics only.

\begin{table}[tbp]
  \centering
  \caption{The three implementations. MLS: moving least squares; TVD: total variation diminishing; SUPG: streamline-upwind Petrov--Galerkin.}
  \label{tab:codes}
  \small
  \setlength{\tabcolsep}{3pt}
  \begin{tabular}{>{\raggedright\arraybackslash}p{2.1cm}>{\raggedright\arraybackslash}p{3.6cm}>{\raggedright\arraybackslash}p{4.0cm}>{\raggedright\arraybackslash}p{4.0cm}}
    \toprule
     & Python & ANSYS (user element) & Abaqus (user element) \\
    \midrule
    $\phiD$, Eq.~\eqref{eq:34} & finite volumes, RK4 & neighboured element method (MLS operators on the Gauss points, 30 neighbours), explicit Euler & Galerkin hexahedra (8 nodes, $2\times2\times2$ Gauss points), implicit Euler, Newton \\
    $c$, $a$, Eq.~\eqref{eq:35} & 7-point Laplacian, sparse direct & MLS Laplacian, direct solver & Galerkin, same Newton system \\
    mechanics & hexahedra, selective reduced integration & hexahedra with the growth material routine & $\bar{\mathbf B}$ hexahedra in the same element \\
    front term & upwind (TVD, minmod) & central MLS gradient with stabilisation (Section~\ref{sec:front}) & Galerkin with SUPG \\
    point model & \multicolumn{3}{l}{the published model \citep{Klempt2026ContinuumBacterialGrowth}, the same Fortran routine in all three} \\
    \bottomrule
  \end{tabular}
\end{table}

\subsection{Code verification}
\label{sec:codever}

\paragraph{Point model in the field.} With uniform fields the coupled run reproduces the point model alone to $2\times10^{-14}$ (Python, two and four species) and $1.8\times10^{-12}$ (ANSYS); the hand-over of $\lv$ to the consumption is exact.

\paragraph{Manufactured solution.} Sources for $\phiD$ and $c$ were built from a manufactured solution \citep{Roache2002MMS} by automatic differentiation and checked symbolically (agreement $10^{-14}$). The finite-volume and the Galerkin code converge at second order in $L_2$ (Figure~\ref{fig:mms}): observed orders between $8^3$, $16^3$ and $32^3$ of $1.85$ and $1.90$ for $\phiD$ and $2.03$ and $2.01$ for $c$ in Python, $1.68$ and $1.81$ for $\phiD$ and $2.00$ and $2.05$ for $c$ in Abaqus; on $32^3$ the errors of $\phiD$ are $2.0\times10^{-4}$ and $2.9\times10^{-4}$. For the SUPG form the order of $\phiD$ is about $1.5$: the largest error sits in a thin layer next to the wall opposite the supply, where $\nabla c\to\mathbf 0$ and the front direction changes over a few cells. The analytic tangent of the Abaqus element agrees with central differences to $5\times10^{-11}$. The neighboured element does not converge next to the walls with the manufactured solution (Appendix~\ref{app:element}).

\begin{figure}[tbp]
  \centering
  \includegraphics[width=\linewidth]{\figs fig_mms.pdf}
  \caption{Manufactured solution: $L_2$ errors of (a)~$\phiD$ and (b)~$c$ against the cell size $h$ for the finite-volume (Python) and the Galerkin (Abaqus) code, with a line of slope 2.}
  \label{fig:mms}
\end{figure}

\paragraph{Mechanics alone.} The last growth field of an Abaqus run, solved by the Python mechanics, gives the seed von Mises stress to $1.1\,\%$ (no obstacles) and $0.5\,\%$ (obstacles with frictionless contact). Differences in the stress between the codes therefore come from the growth field, not from the mechanics.

\subsection{The front term}
\label{sec:front}

The front term of Eq.~\eqref{eq:34} is an advection with the speed $r\,c/(k+c)$ along $\nc$ against a small diffusion: the cell P\'eclet number $r h/(2\beta)$ is about $31$ on a $16^3$ grid ($h=2$~mm$/n$), and a central discretisation oscillates. Four treatments were compared, none of which changes Eqs.~\eqref{eq:34} to \eqref{eq:36}:
\begin{enumerate}
\item a constant added diffusion $r h/4$ (first order; the element as delivered needs it to stay stable);
\item TVD upwinding (MUSCL with the minmod limiter; Python);
\item SUPG \citep{BrooksHughes1982SUPG} with $\tau=\tfrac{h}{2|v|}(\coth Pe-1/Pe)$ on the strong residual (Abaqus);
\item for the explicit element, either the Lax--Wendroff term $\tfrac{\Delta t}{2}\,v\cdot\nabla(v\cdot\nabla\phiD)$ with an added diffusion switched on by an extremum sensor, or upwind-biased MLS gradients.
\end{enumerate}
Figure~\ref{fig:convergence} shows the seed means of $\alpha-1$ and of the von Mises stress at $T^*=1$ of the separate example with antibiotic (Section~\ref{sec:setup}) against the grid for every scheme and code. The growth variable converges for every scheme, and the codes agree within their grid-convergence bands: Python extrapolates to $9.0$ to $9.8\times10^{-4}$, Abaqus to $8.8$ and $9.45\times10^{-4}$, the element to $8.7\times10^{-4}$ with upwind gradients. The stress does not converge with a constant added diffusion: $r h/4$ smears the edge of the seed over a few cells, and the stress, which is set by the jump of $\alpha$ across that edge, follows the smearing; on $16^3$ it is about a quarter of the converged value ($1.3$ to $1.7$ against $5.9$ to $6.3\times10^{-4}$~Pa from the two second-order schemes of two codes), and on $96^3$ it is still rising. Treatments 2 to 4 give the same stress within a few per cent from $32^3$ on (Python TVD $5.90$ and $5.98\times10^{-4}$~Pa on $32^3$ and $64^3$, Abaqus SUPG $6.27$ and $6.22\times10^{-4}$~Pa on $32^3$ and $48^3$, GCI $1.2\,\%$). The element with upwind gradients gives $2.14$, $4.56$ and $5.15\times10^{-4}$~Pa on $8^3$, $16^3$ and $24^3$, extrapolated $5.8\times10^{-4}$~Pa (GCI $16.5\,\%$), within the band of the second-order schemes; its MLS operators over 30 neighbours give it a resolution of about one element, not one Gauss-point spacing (with the same added diffusion, the element on $16^3$ and Python on $32^3$ differ by $1\,\%$ in $\alpha-1$ and $5\,\%$ in the stress). For the stress the treatment of the front term matters more than the grid.

\begin{figure}[tbp]
  \centering
  \includegraphics[width=\linewidth]{\figs fig_convergence.pdf}
  \caption{Separate example with antibiotic: seed means of (a)~$\alpha-1$ and (b)~the von Mises stress at $T^*=1$ against the grid, one line per scheme and code (dashed: constant added diffusion $r h/4$; solid and dotted: the second-order and the stabilised explicit forms). Grey band: the finest grids of the two second-order schemes. These runs start newly reached points from the seed composition (the alternative rule of Section~\ref{sec:newpoints}) at a field step of $0.00125$.}
  \label{fig:convergence}
\end{figure}

\subsection{Solution verification}
\label{sec:solver}

\paragraph{Time step.} In the element codes the point model advances in every increment ($0.00125$ in $T^*$; Abaqus: observed order $1.1$ in time, error against the extrapolated value $0.8\,\%$ in $\alpha-1$ and $2.2\,\%$ in the stress). In the Python code the point model and $\lv$ are updated once per outer step $\Delta T^*$, inside which $\phiD$, $\alpha$ and the fields take explicit RK4 steps. The composition and $\alpha-1$ change by less than $3$ and $0.3\,\%$ between outer steps of $0.02$ and $0.0025$; the stress does not: on $64^3$ the separate example with antibiotic gives $5.68$, $5.82$, $5.94$ and $6.00\times10^{-4}$~Pa at $\Delta T^*=0.025$, $0.01$, $0.005$ and $0.0025$ (first-order extrapolation $6.07$), the mixed one $3.23$, $3.09$, $2.90$ and $2.82\times10^{-4}$~Pa ($2.74$). The sensitivity comes from the front speed: after the antibiotic onset $\lv$ is small, $\nabla c$ is weak and the front direction responds to small changes of the consumption (Appendix~\ref{app:timestep}). The grid series of the examples are therefore run at $\Delta T^*=0.0025$.

\paragraph{Grid.} Table~\ref{tab:gci} gives the grid series of all two-species examples at $\Delta T^*=0.0025$ with the three-grid index of \citet{Celik2008GCI} (safety factor $1.25$), and Figure~\ref{fig:grid} the same data against the grid. The growth converges with a GCI of $0.7$ to $3.3\,\%$ in all examples. The seed stress converges in the separate example, with a GCI of $0.7\,\%$ without and $8.3\,\%$ with the antibiotic. In the mixed examples it converges slowly, with an observed order of about $0.75$ and a GCI of $26$ to $29\,\%$; with the start fractions of \citet{Fritsch2025BayesianMicrofilms} the $128^3$ grid is not yet in the asymptotic range. There the front carries intermediate values of $\phiD$ at a nearly discontinuous edge, and the largest $\phiD$ in the seed still rises with the grid ($0.60$, $0.69$, $0.73$). No run on $32^3$ to $128^3$ left $[0,1]$ by more than $2\times10^{-4}$.
\TODO{Fourth grid ($192^3$, step limit of Appendix~\ref{app:timestep}) for the mixed example with the start fractions of Fritsch et al.; add to Table~\ref{tab:gci} and Figure~\ref{fig:grid} when available.}

\begin{table}[tbp]
  \centering
  \caption{Grid convergence (Python, TVD, outer step $0.0025$): seed means at $T^*=1$, observed order $p$, extrapolated value and GCI \citep{Celik2008GCI}. Mixed: equal amounts unless the start fractions $(\phi_1,\phi_2)$ are given; antib.: with antibiotic; --: not yet in the asymptotic range.}
  \label{tab:gci}
  \small
  \setlength{\tabcolsep}{3.5pt}
  \begin{tabular}{llcccccc}
    \toprule
    example & quantity & $32^3$ & $64^3$ & $128^3$ & $p$ & extrapolated & GCI \\
    \midrule
    separate, antib. & $\alpha-1$ [$10^{-4}$] & $8.453$ & $8.931$ & $9.112$ & $1.4$ & $9.22$ & $1.5\,\%$ \\
    separate, antib. & von Mises [$10^{-4}$~Pa] & $5.523$ & $6.000$ & $6.284$ & $0.75$ & $6.70$ & $8.3\,\%$ \\
    separate & $\alpha-1$ [$10^{-4}$] & $8.729$ & $9.014$ & $9.110$ & $1.6$ & $9.16$ & $0.7\,\%$ \\
    separate & von Mises [$10^{-4}$~Pa] & $7.440$ & $6.730$ & $6.585$ & $2.3$ & $6.55$ & $0.7\,\%$ \\
    mixed, antib. & $\alpha-1$ [$10^{-4}$] & $7.42$ & $8.04$ & $8.32$ & $1.2$ & $8.54$ & $3.3\,\%$ \\
    mixed, antib. & von Mises [$10^{-4}$~Pa] & $2.00$ & $2.82$ & $3.30$ & $0.76$ & $3.99$ & $26\,\%$ \\
    mixed & $\alpha-1$ [$10^{-4}$] & $7.46$ & $8.08$ & $8.35$ & $1.2$ & $8.56$ & $3.1\,\%$ \\
    mixed & von Mises [$10^{-4}$~Pa] & $2.13$ & $3.08$ & $3.65$ & $0.74$ & $4.50$ & $29\,\%$ \\
    mixed $(0.25,0.30)$, antib. & $\alpha-1$ [$10^{-4}$] & $7.42$ & $8.04$ & $8.31$ & $1.2$ & $8.51$ & $3.1\,\%$ \\
    mixed $(0.25,0.30)$, antib. & von Mises [$10^{-4}$~Pa] & $2.01$ & $2.96$ & $4.01$ & -- & -- & -- \\
    \bottomrule
  \end{tabular}
\end{table}

\begin{figure}[tbp]
  \centering
  \includegraphics[width=\linewidth]{\figs fig_grid_examples.pdf}
  \caption{Grid convergence of the two-species examples (Table~\ref{tab:gci}): seed means of (a)~$\alpha-1$ and (b)~the von Mises stress at $T^*=1$ against the grid. Dashed: Richardson extrapolation where the three grids are in the asymptotic range.}
  \label{fig:grid}
\end{figure}

\section{Results}
\label{sec:results}

\subsection{Set-up}
\label{sec:setup}

The examples use a 2~mm cube with a cubic seed of $\phiD=1$ at its centre, the nutrient held at $c=1$ on the face $y=-1$~mm and the antibiotic held at $a=1$ on all outer faces from $T^*=1/3$. In the separate example species~A fills the seed and species~B four small blocks along the edges of the cube; in the mixed example both start in the seed, in equal amounts ($\phi_1=\phi_2=0.275$) or in the proportions of \citet{Fritsch2025BayesianMicrofilms} ($0.25$, $0.30$). The front term is on in every run. The values below are seed means at $T^*=1$ from the Python code (TVD, newly reached points inherit the living composition of their neighbours), which reaches the finest grids; the element gives the same composition within its grid error and its growth within $4\,\%$ after extrapolation. The stress is solved on $16^3$ hexahedra with the element's integration; solving the same $64^3$ field on $32^3$ and $64^3$ elements changes the seed von Mises stress by $0.2\,\%$.

\subsection{Composition and living biomass}
\label{sec:composition}

Table~\ref{tab:examples} gives the seed means and Figure~\ref{fig:panels} the fields of the separate and the mixed example with antibiotic. Without antibiotic the living volume fraction of the seed stays near $0.96$; the antibiotic lowers it to $0.15$ when the species start apart and to $0.04$ when they start mixed. Species~B, the more sensitive one ($b_2=2$), is removed first. The nutrient field decides where the antibiotic acts: the point model compares $c^*=c^*_0c$ with $\alpha^*b_i$, so where the nutrient is high, A survives, and where it is low, A dies as well. When the biomass dies the nutrient is no longer consumed and rises in the seed, from $0.54$ to $0.91$ in the separate example. Living volume fraction, shares and nutrient change by less than $3\,\%$ between the grids ($32^3$ to $128^3$) and the outer steps ($0.02$ to $0.0025$) used.

\begin{table}[tbp]
  \centering
  \caption{Seed means at $T^*=1$, Python, $64^3$ cells, outer step $0.025$; converged growth and stress in Table~\ref{tab:gci}.}
  \label{tab:examples}
  \small
  \setlength{\tabcolsep}{4pt}
  \begin{tabular}{lccccc}
    \toprule
    example & $\lv$ & $\phi_1/(\phi_1+\phi_2)$ & $c$ & $\alpha-1$ & von Mises [Pa] \\
    \midrule
    separate & $0.965$ & $0.978$ & $0.536$ & $9.06\times10^{-4}$ & $6.56\times10^{-4}$ \\
    separate, antibiotic & $0.152$ & $0.951$ & $0.911$ & $8.95\times10^{-4}$ & $5.68\times10^{-4}$ \\
    mixed, equal amounts & $0.961$ & $0.859$ & $0.608$ & $8.33\times10^{-4}$ & $3.50\times10^{-4}$ \\
    mixed, equal amounts, antibiotic & $0.038$ & $0.610$ & $0.973$ & $8.30\times10^{-4}$ & $3.23\times10^{-4}$ \\
    mixed $(0.25,0.30)$, antibiotic & $0.037$ & $0.510$ & $0.973$ & $8.30\times10^{-4}$ & $3.22\times10^{-4}$ \\
    \bottomrule
  \end{tabular}
\end{table}

\begin{figure}[p]
  \centering
  \includegraphics[width=\linewidth]{\assets fig_coupled_3d_panels_separate_ab.pdf}\\[2mm]
  \includegraphics[width=\linewidth]{\assets fig_coupled_3d_panels_mixed_ab.pdf}
  \caption{Separate (upper half) and mixed (lower half) example with antibiotic, Python, $64^3$, one octant removed, at $T^*=0.325$ (first row of each half) and $T^*=1$ (second row): (a)~share $\phi_1/(\phi_1+\phi_2)$, (b)~living volume fraction $\lv$, (c)~nutrient $c$, (d)~antibiotic $a$, (e)~von Mises stress. In the mixed example the living volume fraction falls almost everywhere after the antibiotic.}
  \label{fig:panels}
\end{figure}

\subsection{Antibiotic over time}
\label{sec:schedule}

With the same total dose given constantly, in pulses or as a sine from $T^*=1/3$ (Figure~\ref{fig:schedule}; separate example, $64^3$), $\alpha-1$ changes by $0.1\,\%$ and the seed stress by $0.8\,\%$ between the three schedules, while the living volume fraction of the seed follows the dose: in the pauses of the pulsed dose it recovers from $0.17$ to up to $0.34$ and ends at $0.30$, and the sine gives a lower minimum ($0.11$ against $0.15$) and a higher end value ($0.19$) than the constant dose ($0.15$). How the dose is given acts on the composition and the living biomass, hardly on the mechanics.

\begin{figure}[tbp]
  \centering
  \figorbox{\figs fig_ab_schedule.pdf}{width=\linewidth}
  \caption{The antibiotic given in three ways with the same total dose from $T^*=1/3$ (separate example, $64^3$): (a)~held value at the outer faces, (b)~living volume fraction $\lv$ of the seed over time, (c)~seed means of $\alpha-1$ and of the von Mises stress at $T^*=1$ (each axis spans $\pm5\,\%$ of its mean).}
  \label{fig:schedule}
\end{figure}

\subsection{Which composition does new biofilm receive?}
\label{sec:newpoints}

The rule for newly reached points (Section~\ref{sec:coupling}, item~4) is not fixed by the papers. With the start values of the seed instead of the living composition of the neighbours, the seed values of the cube examples change by a few per cent ($\lv$ $0.147$ against $0.152$ in the separate example with antibiotic). In the maze of \citet[Fig.~11]{Klempt2024DiffusionDrivenGrowth}, where the biofilm keeps spreading after the antibiotic, the rule decides the outcome (Figure~\ref{fig:mazerules}, set-up in Section~\ref{sec:geometry}): the living volume fraction of the colonised region falls to $0.10$ shortly after the onset of the antibiotic at $T^*=20$ and recovers to $0.35$ at $T^*=60$ when new cells take the state of their neighbours, because the biofilm keeps spreading with living cells, and it stays at $0.04$ when they start from the start state of the colony. The Galerkin element with the start-state rule gives the same (0.04 to 0.07 in both codes). Without antibiotic the two codes agree in the colonised area within $3$ to $5\,\%$ and in the living fraction within $0.006$.

\begin{figure}[tbp]
  \centering
  \figorbox{\figs fig_maze_rules.pdf}{width=\linewidth}
  \caption{The maze with the antibiotic held at the entrance from $T^*=20$, at $T^*=60$: living volume fraction $\lv$ where $\phiD\ge0.05$ when newly reached cells (a)~take the state of their colonised neighbours and (b)~take the start state of the colony; walls grey, nutrient source outlined in blue, antibiotic source in red. (c)~Living volume fraction averaged over the colonised cells against $T^*$ for both rules; the antibiotic starts at the dotted line.}
  \label{fig:mazerules}
\end{figure}

\subsection{What reaches the stress}
\label{sec:stress}

The growth $\alpha-1$ hardly depends on the composition: Eq.~\eqref{eq:36} is unchanged, and the antibiotic lowers $\psi_i$ but not $\phiD$ (Table~\ref{tab:gci}: $9.16$ against $9.22\times10^{-4}$ in the separate example, about $1\,\%$). The composition can reach the stress through the nutrient and the front term, but in these examples the effect of the antibiotic on the converged seed stress is within the numerical uncertainty: $6.55$ without and $6.70\times10^{-4}$~Pa with the antibiotic in the separate example (extrapolated, GCI $0.7$ and $8.3\,\%$). On a fixed grid the difference is larger ($6.73$ against $6.00\times10^{-4}$~Pa on $64^3$) because the two cases converge from opposite sides. The mixed seed carries about half the stress of the separate one, because its edge is less sharp: the front carries intermediate values of $\phiD$. A limitation follows from Eq.~\eqref{eq:34}: the front speed $r\,c/(k+c)$ depends on the nutrient only, so dead biofilm keeps spreading after the antibiotic; in the maze about two thirds of the cells reached after the antibiotic hold only dead biofilm when new cells start from the start state. Tying the front speed to the living biomass changes the model and is left as outlook (Section~\ref{sec:discussion}).

\subsection{Geometry: rigid bars and a maze}
\label{sec:geometry}

\citet{Klempt2024DiffusionDrivenGrowth} also grow the biofilm against rigid obstacles (their Section~4.3) and through a maze (Section~4.4). Both were repeated with the coupled model in the Python code, without antibiotic unless stated (Figure~\ref{fig:geometry}).

\paragraph{Rigid bars.} Four bars after Fig.~9 of the paper, with positions read off the sketch: two along $x$ and two along $z$, each $0.25$~mm square, between $0.5$ and $0.75$~mm from the centre. They hold no biofilm and no nutrient, and the biofilm slides along them without friction. All other settings are those of Section~\ref{sec:setup} ($32^3$, outer step $0.0025$). The biofilm grows around the bars and closes behind them, and by $T^*=0.2$ it fills the cube between them. The composition of the seed hardly changes. The bars restrain the swelling: the seed von Mises stress is $7.1\times10^{-3}$~Pa against $7.4\times10^{-4}$~Pa without bars in the separate example, about ten times higher, and $9.5\times10^{-4}$ against $2.1\times10^{-4}$~Pa in the mixed one; the growth $\alpha-1$ of the seed is $9$ and $21\,\%$ lower. These are values on one grid; on $32^3$ the stress without bars lies $14\,\%$ above its extrapolated value (Table~\ref{tab:gci}), small against the factor caused by the bars.

\paragraph{Maze.} The maze of Fig.~11 of the paper, read off the figure ($560~\mu$m across), in two dimensions with $4~\mu$m cells, the constants of Table~6 of the paper ($\beta=10~\mu$m$^2/T^*$, $r=200~\mu$m$/T^*$, $k_\alpha=10^{-6}\,T^{*-1}$), the nutrient held in the centre, the colony of both species at the entrance and the point model of the two-species examples. With the antibiotic, $a=1$ is held at the entrance from $T^*=20$. The model is two-dimensional, so it gives no stress. The biofilm follows the corridors, also into dead ends, and approaches the centre: at $T^*=60$ the nearest point with $\phiD\ge0.5$ is $42~\mu$m from the centre without antibiotic and $133~\mu$m with it. The living volume fraction of the colonised region is $0.92$ without antibiotic; with it, the rule for new cells decides whether the biofilm recovers (Section~\ref{sec:newpoints}).

\begin{figure}[p]
  \centering
  \includegraphics[width=0.82\linewidth]{\assets fig_obstacle_columns_3d.pdf}\\[2mm]
  \includegraphics[width=\linewidth]{\assets fig_maze_k24_h4.pdf}
  \caption{Upper half: growth against four rigid bars (dark grey) after Fig.~9 of \citet{Klempt2024DiffusionDrivenGrowth}, Python, $32^3$, no antibiotic. (a)~Set-up: start biomass (red) and nutrient supply on the face $y=-1$~mm (green). (b)--(d)~Surface $\phiD=0.5$ (separate example) and $0.2$ (mixed example, which starts at $\phiD=0.55$) at $T^*=0.05$, $0.1$ and $0.2$. Lower half: the maze of Fig.~11 of the same paper with the coupled model (two dimensions, $4~\mu$m cells): $\phiD$ at $T^*=12$, $24$, $36$ and $60$, (a)~without and (b)~with the antibiotic at the entrance from $T^*=20$. Walls white, nutrient source outlined in red.}
  \label{fig:geometry}
\end{figure}

\subsection{A clinical molar crown}
\label{sec:realtooth}

The cube examples place the seed in a box. As a first step towards the geometry of the application, the coupled model was also run on the crown of a human molar from the Open-Full-Jaw dataset \citep{Gholamalizadeh2022OpenFullJaw}, with the biofilm seeded on the occlusal surface and the nutrient supplied from the surrounding volume (Figure~\ref{fig:realtooth}). The constants are those of Table~\ref{tab:parameters}, and the two species start as in the separate example. The run shows the biofilm following the fissures of the crown, and the living volume fraction, the share and the stress at the end of the antibiotic.
\TODO{Section~\ref{sec:realtooth}: run on the molar crown in progress (another script); fill in the set-up (mesh size, seed and sources, time), the seed and fissure values of $\lv$, $\phi_1/(\phi_1+\phi_2)$ and the stress, and check the wording against the figure. Dataset licence CC BY-NC-SA: state it in the data availability section.}

\begin{figure}[tbp]
  \centering
  \figorbox{\figs fig_real_tooth_coupled.pdf}{width=\linewidth}
  \caption{The coupled model on the crown of a human molar \citep{Gholamalizadeh2022OpenFullJaw}: biofilm $\phiD$, living volume fraction $\lv$, share $\phi_1/(\phi_1+\phi_2)$ and von Mises stress at $T^*=1$ with the antibiotic from $T^*=1/3$. Caption to be completed with the figure.}
  \label{fig:realtooth}
\end{figure}

\section{Discussion}
\label{sec:discussion}

\paragraph{What the coupling shows.} The composition reaches the growth field only through the consumption $g\,\phiD(\lv)\,c$. When the antibiotic kills the biomass, less nutrient is consumed, the nutrient rises in the seed and the front term moves the edge differently. The effect on the growth is small, about $1\,\%$, because the growth law $\dot\alpha=k_\alpha\phiD$ counts all biomass, living or dead; the effect on the seed stress is about $10\,\%$, because the stress is set by the jump of $\alpha$ across the edge of the seed, which the front shapes. The local nutrient decides where the antibiotic acts: the same dose removes a species in one part of the biofilm and not in another, which a model without space cannot show. The stress depends on how sharp the edge of the biofilm is, and that edge is set by the front term, so the coupling needs a second-order treatment of that term; a constant added diffusion gives a quarter of the converged seed stress on $16^3$ and does not converge up to $96^3$.

\paragraph{Limitations.} (i)~The front speed depends on the nutrient only, so dead biofilm keeps spreading after the antibiotic. Tying the speed to the living biomass is a change of the model; in a test on $16^3$ (front speed multiplied by the local living fraction) the seed stress rose by $51\,\%$ in the separate example and about fourfold in the mixed one, so this choice matters for the stress. (ii)~Eq.~\eqref{eq:36} lets dead biomass grow at the same rate as living biomass; growth from the living biomass, as in \citet{WannerGujer1986}, would make the antibiotic act on the stress directly. (iii)~The point model has no space; the composition moves only through the rule for newly reached points, which decides in the maze whether the biofilm recovers. A species-resolved continuum would replace the rule. (iv)~The mechanics is coupled one way, as in \citet{Klempt2024DiffusionDrivenGrowth}; with $E=10$~Pa the neglected term acts only for coefficients far beyond the paper values (an effect on the growth from about $10^5$ on, $16^3$ test). (v)~Living and dead biomass, and both species, have the same stiffness; a species-dependent stiffness changes the seed stress by $9$ to $89\,\%$ in a $16^3$ test, depending on which species is stiffer. (vi)~The examples use synthetic data with known parameters, the stress converges slowly when the species start mixed, and the element's values carry a grid error of $16.5\,\%$. Calibrated data \citep[for example the five-species data of][]{Heine2025PeriImplant} and a measured stress are needed before the stresses can be validated, not only verified.

\paragraph{Outlook.} Species-resolved transport, so that the composition moves with the biomass; front speed and growth tied to the living biomass, with the consequences for the stress quantified above; growth laws with a mechanical feedback; calibrated multi-species data in the element; and a realistic tooth or implant geometry. Each can be added to the coupling as presented here, and the verification procedure, several codes, a manufactured solution and grid convergence indices, carries over.

\section{Conclusions}
\label{sec:conclusions}

This work coupled the multi-species point model of \citet{Klempt2026ContinuumBacterialGrowth} to the continuum growth model of \citet{Klempt2024DiffusionDrivenGrowth} at every integration point, without changing the equations of either model. The point model receives the local nutrient and antibiotic from the fields, $c^*=c^*_0c$ and $\alpha^*=\alpha^*_0a$, and its living volume fraction sets the uptake $g\,\phiD(\lv)\,c$ of both.
\begin{enumerate}
\item The coupling was implemented in a finite-volume code, in the neighboured-element user element of the original authors' group in ANSYS and in a Galerkin user element in Abaqus. The point model inside each code reproduces the point model alone to round-off in uniform fields, the finite-volume and the Galerkin code converge at second order for a manufactured solution, and the three codes agree within their grid-convergence bands. Three numerical details of the neighboured element had to be corrected for this.
\item The front term \fterm{} sets the accuracy of the stress. A constant added diffusion keeps it stable but gives a quarter of the converged seed stress; a second-order upwind treatment is needed. With it, the growth converges with a GCI of $0.7$ to $3.3\,\%$ in all examples, the seed stress with $0.7$ to $8.3\,\%$ when the species start apart and $26$ to $29\,\%$ when they start mixed.
\item An antibiotic removes most of the living biomass of the seed (from $0.97$ to $0.15$ when the species start apart, from $0.96$ to $0.04$ when mixed), while the growth changes by about $1\,\%$ and the seed stress by about $10\,\%$: the antibiotic lowers the consumption, more nutrient reaches the front, and the edge of the seed changes. The local nutrient decides where the antibiotic acts, because the point model weighs $c^*$ against $\alpha^*b_i$.
\item The composition of biofilm that reaches a new point is not fixed by either model. Taking the living composition of the neighbours or the start state gives the same cube results within a few per cent, but decides in a maze whether the biofilm recovers after an antibiotic (living volume fraction $0.35$ against $0.04$).
\item Rigid obstacles restrain the swelling and raise the seed stress about tenfold; a maze makes the biofilm follow its corridors and approach the nutrient source.
\end{enumerate}
The composition reaches growth and stress only through the nutrient and the front speed, because the growth law counts all biomass. The next steps are models in which the front speed, the growth or the stiffness depend on the living biomass or on the species.

\section*{Data availability}
The finite-volume code, the point-model routine, the run scripts, the result files of the examples and the figure scripts are in the public repository \url{https://github.com/keisuke58/pde-fem-biofilm}. The examples use the published synthetic data of \citet{Fritsch2025BayesianMicrofilms}; no calibration data are used in this work. The user elements depend on the commercial codes and on the element of \citet{Rudolf2025NEM}; the material routines are in the repository.
\TODO{Confirm with all authors which parts of the element code may be published, and whether a release (DOI) of the repository is made at submission.}

\section*{Declarations}
Author list, contributions, funding and conflicts of interest: to be agreed.
\TODO{Authors, affiliations, contributions, funding (LUH/IKM, Keio), acknowledgements; apply the Springer Nature template (sn-jnl.cls) of Biomechanics and Modeling in Mechanobiology at submission.}

\appendix

\section{Numerical details}
\label{app:numerics}

\subsection{The neighboured element}
\label{app:element}

Three details of the element were corrected, without changing its equations: the surface points take the value of the nearest Gauss point instead of a distance-weighted mean, the Laplacian is used in a conservative form, and the neighbours of the fit are chosen by distance instead of by element number. Before, $\phiD$ leaked through the walls of the cube (the surface points took an inverse-square mean, and the column sums of the Laplacian were not zero), and the stencils of the $16^3$ and $24^3$ grids were asymmetric inside the domain. With a manufactured solution the corrected element still does not converge from $8^3$ to $16^3$ ($L_2$ error of $\phiD$ $2.2\times10^{-3}$ to $1.8\times10^{-3}$, of $c$ $8.2\times10^{-3}$ to $1.2\times10^{-2}$): the MLS stencils of the Gauss points next to a face include the surface points and are ill-conditioned (error of the Laplacian of the manufactured $c$: $0.88$ on $8^3$, $65$ on $16^3$, against $|\nabla^2c|=3.5$). Rebuilding these stencils with the neighbours mirrored across the faces (zero flux built in) makes the operators consistent up to the walls, and with the weight width proportional to $h$ the element converges ($\phiD$ $9.5\times10^{-4}$ to $3.4\times10^{-4}$, order $1.5$; $c$ $9.0\times10^{-4}$ to $2.2\times10^{-4}$, order $2.0$). In the coupled example the mirrored operators run on $8^3$ but become unstable at a corner of the $16^3$ grid after the antibiotic onset; they are not used for the results. The explicit update of the element is limited by $\beta\Delta t/h^2\le0.12$ to $0.16$ on $16^3$ and $24^3$ (Bloch analysis of the stencil); the Lax--Wendroff form with the extremum sensor keeps $\phiD$ in $[-0.03,1.01]$ on $64^3$ in Python but is unstable at the corners of the supply face on $24^3$ in the element, so the element's stresses are quoted from the upwind MLS gradients. The time-step series of the element and the stability analysis will be given in the supplementary material.

\subsection{Outer step, regularisation and stability}
\label{app:timestep}

Within an outer step of the Python code, the inner RK4 step, a trapezoidal update of $\lv$ over the step and letting newly reached points consume at once each change the stress by less than $1\,\%$; the one-increment lag of $\lv$ and the regularisation of the front direction ($\varepsilon^2=10^{-6}$ against $10^{-8}$, $0.1\,\%$) are not the cause of the step sensitivity of Section~\ref{sec:solver}. A larger regularisation ($\varepsilon^2=10^{-4}$) changes the mixed stress by $30\,\%$, which shows that the front direction, not the speed, carries the sensitivity. Without antibiotic the outer step changes the stress by $2.5\,\%$ (separate) and $12\,\%$ (mixed). The explicit inner step is limited by the diffusion number $D=\beta\Delta t/h^2\le0.15$ and the Courant number of the front $Co\le0.4$; on $192^3$ both limits are close at the same time and the step diverged in the first outer step, which the combined condition $12D+2Co\le2.2$ removes without changing the step on $32^3$ to $128^3$. The regularised front direction and the P\'eclet number of the examples are unchanged by these choices.

\section*{Supplementary material}
The time-step series of the element and of the Galerkin code, the stability analysis of the explicit element, the sensitivities to $\phi_{\min}$, $d$ and the start rule, and the animation of the maze will be provided as supplementary material.

{\small
\bibliographystyle{plainnat}
\bibliography{refs}
}

\end{document}
